% GENERATED FILE. Edit canonical lessons, then run npm run generate:books.
% canonical-source-hash: fnv1a64:041cd5aa27dbc22a
% canonical-lessons: GE-C471-traene, GE-C471-schluchzen, GE-C471-trost, GE-C471-troesten, GE-C471-mitleid

\chapter{Tränen und Trost}
\label{ch:ge-tranen-und-trost}
\hlchaptermodality{471}

\begin{chapteropening}
\textbf{By the end of this chapter:} \emph{I can talk about crying and comforting someone.}

Complete the last lesson of chapter 471: I can talk about crying and comforting someone.
\end{chapteropening}

\section[die Träne]{die Träne — tear (from crying)}

\label{lesson:GE-C471-traene}



\begin{quote}
Before the new one: say the German for excitement, then the German for to fidget.
\end{quote}

\subsection*{You'll want to know: die Träne}
\textbf{die Träne} — \textquotedblleft{}tear (from crying)\textquotedblright{}.

\begin{quote}
Beim Abschied hatte sie Tränen in den Augen.
\end{quote}

\begin{quote}
\emph{When they said goodbye, she had tears in her eyes.}
\end{quote}

Plural: \emph{die Tränen}. \emph{In Tränen ausbrechen} — to burst into tears.

\begin{cousinweb}
\emph{Träne} is inherited from Middle High German \emph{træne} \textquotedblleft{}tears\textquotedblright{}, a form of \emph{trahen}, from Old High German \emph{trahan} \textquotedblleft{}tear\textquotedblright{} (EtymDB); it is usually traced to Proto-Germanic \emph{trahnuz} \textquotedblleft{}tear\textquotedblright{}.
\end{cousinweb}

\begin{grammarlens}[title={What you've built}]
One more word for telling, explaining and describing at B1.
\end{grammarlens}

\subsection*{Your turn}
\begin{itemize}
\raggedright
  \item \emph{Say it:} \emph{die Träne}
  \item \emph{Say it:} \emph{die Träne}, once more
  \item \emph{Say it:} a sentence of your own with \emph{die Träne}
  \item \emph{Recall:} say the German for excitement, then the German for to fidget, then say \emph{die Träne} again
\end{itemize}

\begin{tcolorbox}[breakable,colback=teal!4,colframe=teal!35!black,title={Before you move on}]
What does die Träne mean? (\textquotedblleft{}Tear (from crying)\textquotedblright{}.) Say it once more.
\end{tcolorbox}
\section[schluchzen]{schluchzen — to sob}

\label{lesson:GE-C471-schluchzen}



\begin{quote}
Before the new one: say the German for to fidget, then the German for tear.
\end{quote}

\subsection*{You'll want to know: schluchzen}
\textbf{schluchzen} — \textquotedblleft{}to sob\textquotedblright{}.

\begin{quote}
Das Mädchen schluchzte leise in seinem Zimmer.
\end{quote}

\begin{quote}
\emph{The girl was sobbing quietly in her room.}
\end{quote}

Stronger than \emph{weinen}.

\begin{cousinweb}
\emph{schluchzen} is usually explained as an imitative word beside \emph{schlucken} \textquotedblleft{}to swallow\textquotedblright{}; EtymDB lists it without a source, so the details are not settled.
\end{cousinweb}

\begin{grammarlens}[title={What you've built}]
One more word for telling, explaining and describing at B1.
\end{grammarlens}

\subsection*{Your turn}
\begin{itemize}
\raggedright
  \item \emph{Say it:} \emph{schluchzen}
  \item \emph{Say it:} \emph{schluchzen}, once more
  \item \emph{Say it:} a sentence of your own with \emph{schluchzen}
  \item \emph{Recall:} say the German for to fidget, then the German for tear, then say \emph{schluchzen} again
\end{itemize}

\begin{tcolorbox}[breakable,colback=teal!4,colframe=teal!35!black,title={Before you move on}]
What does schluchzen mean? (\textquotedblleft{}To sob\textquotedblright{}.) Say it once more.
\end{tcolorbox}
\section[der Trost]{der Trost — comfort, consolation}

\label{lesson:GE-C471-trost}



\begin{quote}
Before the new one: say the German for tear, then the German for to sob.
\end{quote}

\subsection*{You'll want to know: der Trost}
\textbf{der Trost} — \textquotedblleft{}comfort, consolation\textquotedblright{}.

\begin{quote}
Ein warmer Tee war ein kleiner Trost.
\end{quote}

\begin{quote}
\emph{A warm tea was a small comfort.}
\end{quote}

No plural. \emph{Ein schwacher Trost} — cold comfort.

\begin{cousinweb}
\emph{Trost} is inherited from Old High German \emph{trōst} \textquotedblleft{}consolation\textquotedblright{}, from Proto-Germanic \emph{traustą} (EtymDB), built on an adjective meaning \textquotedblleft{}firm, strong\textquotedblright{}; English \emph{trust} is a cousin.
\end{cousinweb}

\begin{grammarlens}[title={What you've built}]
One more word for telling, explaining and describing at B1.
\end{grammarlens}

\subsection*{Your turn}
\begin{itemize}
\raggedright
  \item \emph{Say it:} \emph{der Trost}
  \item \emph{Say it:} \emph{der Trost}, once more
  \item \emph{Say it:} a sentence of your own with \emph{der Trost}
  \item \emph{Recall:} say the German for tear, then the German for to sob, then say \emph{der Trost} again
\end{itemize}

\begin{tcolorbox}[breakable,colback=teal!4,colframe=teal!35!black,title={Before you move on}]
What does der Trost mean? (\textquotedblleft{}Comfort, consolation\textquotedblright{}.) Say it once more.
\end{tcolorbox}
\section[trösten]{trösten — to comfort, to console}

\label{lesson:GE-C471-troesten}



\begin{quote}
Before the new one: say the German for to sob, then the German for comfort.
\end{quote}

\subsection*{You'll want to know: trösten}
\textbf{trösten} — \textquotedblleft{}to comfort, to console\textquotedblright{}.

\begin{quote}
Die Oma tröstet das Kind nach dem Sturz.
\end{quote}

\begin{quote}
\emph{Grandma comforts the child after the fall.}
\end{quote}

\emph{Sich trösten mit} — to console yourself with.

\begin{cousinweb}
\emph{trösten} is built on \emph{Trost} \textquotedblleft{}comfort\textquotedblright{}, which is inherited from Old High German \emph{trōst}, from Proto-Germanic \emph{traustą} (EtymDB).
\end{cousinweb}

\begin{grammarlens}[title={What you've built}]
One more word for telling, explaining and describing at B1.
\end{grammarlens}

\subsection*{Your turn}
\begin{itemize}
\raggedright
  \item \emph{Say it:} \emph{trösten}
  \item \emph{Say it:} \emph{trösten}, once more
  \item \emph{Say it:} a sentence of your own with \emph{trösten}
  \item \emph{Recall:} say the German for to sob, then the German for comfort, then say \emph{trösten} again
\end{itemize}

\begin{tcolorbox}[breakable,colback=teal!4,colframe=teal!35!black,title={Before you move on}]
What does trösten mean? (\textquotedblleft{}To comfort, to console\textquotedblright{}.) Say it once more.
\end{tcolorbox}
\section[das Mitleid]{das Mitleid — pity, compassion}

\label{lesson:GE-C471-mitleid}



\begin{quote}
Before the new one: say the German for comfort, then the German for to comfort.
\end{quote}

\subsection*{You'll want to know: das Mitleid}
\textbf{das Mitleid} — \textquotedblleft{}pity, compassion\textquotedblright{}.

\begin{quote}
Ich habe Mitleid mit dem kranken Hund.
\end{quote}

\begin{quote}
\emph{I feel sorry for the sick dog.}
\end{quote}

No plural. \emph{Mitleid haben mit} — to feel sorry for.

\begin{cousinweb}
\emph{Mitleid} is built on \emph{mit} \textquotedblleft{}with\textquotedblright{} and \emph{Leid} \textquotedblleft{}sorrow\textquotedblright{}; EtymDB traces it to Middle High German \emph{mitelīden} \textquotedblleft{}to suffer with\textquotedblright{}. It is usually seen as modelled on Church Latin \emph{compassio} \textquotedblleft{}fellow-suffering\textquotedblright{} (Lewis \& Short).
\end{cousinweb}

\begin{grammarlens}[title={What you've built}]
One more word for telling, explaining and describing at B1.
\end{grammarlens}

\subsection*{Your turn}
\begin{itemize}
\raggedright
  \item \emph{Say it:} \emph{das Mitleid}
  \item \emph{Say it:} \emph{das Mitleid}, once more
  \item \emph{Say it:} a sentence of your own with \emph{das Mitleid}
  \item \emph{Recall:} say the German for comfort, then the German for to comfort, then say \emph{das Mitleid} again
\end{itemize}

\begin{tcolorbox}[breakable,colback=teal!4,colframe=teal!35!black,title={Before you move on}]
What does das Mitleid mean? (\textquotedblleft{}Pity, compassion\textquotedblright{}.) Say it once more.
\end{tcolorbox}
